\documentclass[11pt]{article}
\usepackage[a4paper,margin=27mm]{geometry}
\usepackage[T1]{fontenc}
\usepackage{lmodern,microtype}
\usepackage{amsmath,amssymb,amsthm,mathtools}
\usepackage[colorlinks=true,linkcolor=blue,citecolor=blue,urlcolor=blue]{hyperref}
\newtheorem{theorem}{Theorem}[section]
\newtheorem{lemma}[theorem]{Lemma}
\newtheorem{proposition}[theorem]{Proposition}
\newtheorem{corollary}[theorem]{Corollary}
\theoremstyle{remark}\newtheorem{remark}[theorem]{Remark}
\DeclareMathOperator{\Var}{Var}
\DeclareMathOperator{\rank}{rank}
\DeclareMathOperator{\essinf}{ess\,inf}
\newcommand{\E}{\mathbb E}
\newcommand{\Prob}{\mathbb P}
\newcommand{\N}{\mathbb N}
\newcommand{\Q}{\mathbb Q}
\newcommand{\ind}{\mathbf 1}
\newcommand{\fall}[2]{#1^{\underline{#2}}}
\newcommand{\rise}[2]{(#1)^{\overline{#2}}}
\newcommand{\cO}{\mathcal O}
\newcommand{\cF}{\mathcal F}
\title{Limiting Betti-number distributions in preferential-attachment clique complexes}
\author{Kaiwen Jin\\\small\href{mailto:yui.ui.siki@gmail.com}{\texttt{yui.ui.siki@gmail.com}}}
\date{October 14, 2026}
\begin{document}
\maketitle
\begin{abstract}
We study the finite clique complexes of sequential affine preferential-attachment graphs.
For every $q\ge2$, $m\ge2q$ and $\delta>-m$ with
$x=(m+\delta)/(2m+\delta)<1/(2q)$, the $q$th rational Betti number,
divided by its expectation, converges to a nondegenerate random variable.
We identify its law as a normalized sum of products of attachment weights
over induced octahedral subgraphs in an explicit independent-Beta construction.
The limit is positive almost surely, has zero essential infimum and has every
positive moment. In particular, its upper tail is not regularly varying with
any finite positive exponent. The proof compares all octahedral link
contributions with finite-time homology; larger link cores have lower order,
and an earlier cone makes each fixed base contribute to subsequent births.
\end{abstract}

\section{Introduction}\label{sec:introduction}
Preferential attachment produces growing graphs in which vertices of large
degree attract more edges. Their clique complexes record higher-order
connections: a complete subgraph is filled by a simplex, and the Betti numbers
measure the homology that remains. The variability of these numbers across
independently generated graphs is a different question from their average
growth.

In the affine model, each new vertex adds $m$ edges and the attraction of
an old vertex is its current degree plus $\delta$. Write
$x=(m+\delta)/(2m+\delta)$.
For the sequential version considered here, Siu, Samorodnitsky, Yu and
He~\cite{ssyh} proved that the expected $q$th Betti number has order
$T^{1-2qx}$ when $q\ge2$, $m\ge2q$ and $x<1/(2q)$. Siu~\cite{siu}
then obtained matching upper and lower bounds in probability, up to any
diverging multiplicative factor. The latter proof counts common neighbors of
a fixed induced octahedral sphere and controls lower-order killing terms.
These bounds locate the scale but leave the limiting distribution unidentified.
In Section~1.2 of the author preprint~\cite{siu}, mean-normalized convergence
and a power-law tail were suggested by numerical simulations.

We prove the convergence throughout this strict polynomial-growth regime.
The limit is an explicit random series over all induced octahedral bases,
and it has every positive moment. This gives a nonconstant limiting law while
refuting regular variation with a finite positive exponent. The normalization
is the ensemble expectation of the Betti number, as in the distributional
question; no empirical normalization is involved.

Two issues separate a distributional result from the existing order bounds.
The count must include all base spheres, with a summable bound on their large
labels. It must also track whether their link homology actually produces
new homology in the finite complex. We use the independent-Beta description
of the attachment process to control the first issue. For the second,
deleting vertices of small degree reduces a typical link to a single
octahedral sphere. Every fixed base eventually acquires an earlier cone,
after which its one-dimensional link homology maps to zero in the old complex
and contributes one birth. The remaining link cores require an additional
neighbor and have lower order.

Random environments also describe non-Gaussian triangle-count limits in the
same attachment model, as shown recently by Dey and Terlov~\cite{dey-terlov};
their triangle observable counts edge multiplicity.
Here the coefficient additionally retains the induced graph on the base
vertices. Joint weight and edge estimates handle that dependence, including
overlapping bases in every fixed moment. A conditional estimate obtained by
taking one Beta variable close to one shows that the coefficient can be
arbitrarily small with positive probability, proving nondegeneracy.

Section~\ref{sec:model} states the exact model and the limiting law.
Section~\ref{sec:environment} establishes the weight and counting estimates.
Section~\ref{sec:homology} compares the counts with finite homology, and
Section~\ref{sec:tail} proves nondegeneracy and the tail conclusion.
Homology is over $\Q$ throughout. Constants may depend on the fixed model
parameters and, where indicated, on a fixed moment order; they do not depend
on time or vertex labels.

\section{The model and the limiting law}\label{sec:model}
Fix integers $q\ge2$, $m\ge2q$ and a real number $\delta>-m$ such that
\begin{equation}\label{eq:parameters}
 a=m+\delta>0,\qquad c=2m+\delta=m+a,\qquad
 x=\frac ac<\frac1{2q},\qquad k=2q,\qquad \rho=1-kx>0.
\end{equation}
The graph $G_2$ consists of two vertices joined by $m$ parallel edges.
At each time $t\ge3$, vertex $t$ adds $m$ edges, one at a time, to old
vertices. After the first $b-1$ edges, the next endpoint $v<t$ is chosen
with probability
\begin{equation}\label{eq:sequential}
 \frac{d_{t,b-1}(v)+\delta}
 {2m(t-2)+(t-1)\delta+(b-1)},\qquad 1\le b\le m.
\end{equation}
Degrees count multiplicity. There are no self-loops. Write $X_T$ for the
clique complex of the underlying simple graph, and
$\beta_q(T)=\dim_\Q H_q(X_T;\Q)$.

The following construction puts all times on one probability space.
Let $\psi_1=1$ and let
\begin{equation}\label{eq:beta}
 \psi_j\sim\operatorname{Beta}(a,b_j),\qquad
 b_j=c(j-1)-m,\qquad j\ge2,
\end{equation}
be independent; $b_2=a$ and all $b_j$ are positive.
Given $(\psi_j)$, every edge from $t$ chooses its old endpoint independently
with probabilities
\begin{equation}\label{eq:kernel}
 p(v,t)=\psi_v\prod_{v<j<t}(1-\psi_j),\qquad v<t.
\end{equation}
Their sum is one by telescoping. Independent uniform variables on $(0,1)$
can be used to make these choices. At $t=2$, all endpoints are $1$.
Denote this infinite graph by $G$ and use its simple graph when speaking
of induced subgraphs.

\begin{lemma}\label{lem:model}
The restrictions of this construction have exactly the joint sequential
graph law in~\eqref{eq:sequential}.
\end{lemma}
\begin{proof}
This is the P\'olya representation used in~\cite[Definition~5.1 and
Theorem~5.2 of the author preprint]{siu}. We include the verification to fix
the precise variant. At the completion of the birth of vertex $j$, the
total weight of the vertices before $j$ is
\[
 2m(j-2)+(j-1)\delta+m=c(j-1)-m=b_j,
\]
and vertex $j$ has weight $a$. A later endpoint at $j$ contributes $\psi_j$
to the likelihood. An endpoint before $j$ from a later vertex contributes
$1-\psi_j$; other endpoints do not change its factor. Conditional on any
finite endpoint history, the variables belonging to eligible vertices
therefore remain independent, with posterior law
$\operatorname{Beta}(w_j,W_{j-1})$, where $w_j$ is the current weight
of $j$ and $W_i=\sum_{v\le i}w_v$.
This holds during a partial new birth as well. Taking the posterior
expectation in~\eqref{eq:kernel} yields
\[
 \frac{w_v}{W_v}\prod_{v<j<t}\frac{W_{j-1}}{W_j}
 =\frac{w_v}{W_{t-1}}.
\]
The denominator after $b-1$ new edges is precisely the one
in~\eqref{eq:sequential}. Initially $w_1=w_2=a$ and
$\psi_2\sim\operatorname{Beta}(a,a)$, while the $m$ initial edges are
deterministic. This proves the claim for every finite ordered history.
\end{proof}

We will show that the following limits exist and are positive almost surely:
\begin{equation}\label{eq:weights}
 R_v=\lim_{n\to\infty}n^x\prod_{v<j\le n}(1-\psi_j),
 \qquad D_v=\psi_v R_v.
\end{equation}
They describe the fixed-vertex attachment probabilities:
$t^x p(v,t)\to D_v$.
Let $\cO_q$ be the graph $K_{2q}$ with a perfect matching removed.
Its clique complex is the join of $q$ two-point spaces, an octahedral
$(q-1)$-sphere. For every $k$-element set $S\subset\N$, put
\begin{equation}\label{eq:coefficient}
 I_S=\ind\{G[S]\cong\cO_q\},\qquad
 A=\sum_{\substack{S\subset\N\\ |S|=k}} I_S\prod_{v\in S}D_v.
\end{equation}
Here $G[S]$ is the induced simple graph, so $I_S$ imposes the missing
matching as well as the present edges. The series is initially understood
as a nonnegative extended random variable; its finiteness is part of the result.

\begin{theorem}\label{thm:main}
For every parameter triple in~\eqref{eq:parameters}, the weights
in~\eqref{eq:weights} are positive and finite almost surely, and
\[
 A>0\ \text{almost surely},\qquad
 \essinf A=0,\qquad 0<\E A<\infty,\qquad
 \E A^p<\infty\quad (p>0).
\]
In the construction~\eqref{eq:beta}--\eqref{eq:kernel},
\begin{equation}\label{eq:betti-limit}
 \frac{\beta_q(T)}{T^\rho}
 \xrightarrow{L^1}\frac{\fall{m}{k}}{\rho}A,
 \qquad
 \frac{\E\beta_q(T)}{T^\rho}
 \longrightarrow\frac{\fall{m}{k}}{\rho}\E A,
\end{equation}
where $\fall{m}{k}=m(m-1)\cdots(m-k+1)$.
Consequently the mean-normalized Betti number converges in distribution to
\begin{equation}\label{eq:law}
 W_{q,m,\delta}=\frac{A}{\E A},
\end{equation}
a nondegenerate random variable with mean one. Its upper tail does not satisfy
\begin{equation}\label{eq:rv}
 \lim_{z\to\infty}\frac{\Prob(W_{q,m,\delta}>uz)}
 {\Prob(W_{q,m,\delta}>z)}=u^{-\alpha},\qquad u>0,
\end{equation}
for any finite $\alpha>0$.
\end{theorem}

The law in~\eqref{eq:law} is specified by independent Beta variables,
independent endpoint uniforms and the convergent series~\eqref{eq:coefficient}.
All parameters are fixed before $T\to\infty$.

\section{Attachment weights and octahedral counts}\label{sec:environment}
\subsection{Joint moments}
We first obtain bounds uniform in all vertex labels. For a directed simple
graph $F$ whose edges point from larger to smaller labels, let $d_-(v)$
and $d_+(v)$ be its incoming and outgoing degrees, respectively.
The event $F\subset G$ requires its edges and places no restriction on
other edges. The number of vertices and edges in this lemma is fixed.

\begin{lemma}\label{lem:mixed}
All the limits in~\eqref{eq:weights} exist and are positive almost surely.
For nonnegative integers $r_v$ on a finite set $V$ containing the vertices
of $F$,
\begin{equation}\label{eq:mixed}
 \E\left[\ind\{F\subset G\}\prod_{v\in V}D_v^{r_v}\right]
 \le C\prod_{v\in V}
 v^{-(1-x)(r_v+d_-(v))-x d_+(v)}.
\end{equation}
The constant depends only on $m,\delta$, the finite size bounds and
$\sum_v r_v$, not the labels. Extend $r_v,d_\pm(v)$ by zero outside $V$.
For $K\ge2$, define
\[
 \ell_K=r_K+d_-(K),\qquad
 h_K=\sum_{v<K}r_v+\#\{u\to v\in F:v<K<u\}.
\]
The conditional version of~\eqref{eq:mixed}, with $\psi_K=z\in(0,1)$,
holds with the additional factor $C K^{\ell_K}(1-z)^{h_K}$, uniformly in
$K$ and $z$.
\end{lemma}
\begin{proof}
Write $Y_j=-\log(1-\psi_j)$. For $0\le u<1$,
\[
 0\le-\log(1-u)-u\le\frac{u^2}{1-u},\qquad
 [-\log(1-u)]^2\le\frac{u^2}{(1-u)^2}.
\]
For all sufficiently large $j$, beta integrals give
\begin{align*}
 \E\psi_j&=\frac{a}{cj-2m}=\frac{x}{j}+O(j^{-2}),\\
 \E\frac{\psi_j^2}{1-\psi_j}
 &=\frac{a(a+1)}{(a+b_j)(b_j-1)}=O(j^{-2}),\\
 \E\frac{\psi_j^2}{(1-\psi_j)^2}
 &=\frac{a(a+1)}{(b_j-1)(b_j-2)}=O(j^{-2}).
\end{align*}
Thus $\E Y_j=x/j+O(j^{-2})$ and $\Var(Y_j)=O(j^{-2})$.
The finitely many early indices have finite squared logarithmic moments,
since $b_j>0$. The independent series criterion
\cite[Theorem~2.5.6]{durrett} shows that
$\sum_{j\ge2}(Y_j-\E Y_j)$ converges almost surely.
The summable expectation error and the convergence of the harmonic sum
minus $\log n$ now prove~\eqref{eq:weights}. A countable intersection
handles every $v$.

Given the environment, assign distinct edge slots to the required older
neighbors of each vertex of $F$. A union bound gives
\begin{equation}\label{eq:edge-bound}
 \Prob(F\subset G\mid(\psi_j))
 \le m^{|E(F)|}\prod_{u\to v\in F}p(v,u).
\end{equation}
If a required outdegree exceeds $m$, the event is empty. To estimate the
right side multiplied by the weight product, first use
$D_{v,n}=\psi_v n^x\prod_{v<j\le n}(1-\psi_j)$, with $n$ larger than
all labels. Independence reduces the expectation to beta factors
\begin{equation}\label{eq:beta-factor}
 \E[\psi_j^{\ell_j}(1-\psi_j)^{h_j}]
 =\frac{\rise{a}{\ell_j}\rise{b_j}{h_j}}
 {\rise{a+b_j}{\ell_j+h_j}},
\end{equation}
where a rising factorial is indicated by an overline, and
\[
 \ell_j=r_j+d_-(j),\qquad
 h_j=\sum_{v<j}r_v+\#\{u\to v\in F:v<j<u\}.
\]
Both exponents have a fixed bound. At the finitely many sites with
$\ell_j>0$, the factor is at most
$Cj^{-\ell_j}\exp(-xh_j/j)$. At other sites its logarithm is
$-xh_j/j+O(j^{-2})$. These statements follow by taking logarithms of the
finite products in~\eqref{eq:beta-factor}; the constants also cover the
finitely many small indices. Consequently the complete product is bounded by
\begin{align*}
 Cn^{x\sum_vr_v}\prod_vv^{-r_v-d_-(v)}
 \exp\left(-x\sum_{j=2}^n\frac{h_j}{j}\right).
\end{align*}
For each weight, the harmonic sum is $\log(n/v)+O(1/v)$, and for an edge
$u\to v$ it is $\log(u/v)+O(1/v)$. Substitution gives
\[
 Cn^{x\sum_vr_v}\prod_vv^{-r_v-d_-(v)}
 \prod_v(n/v)^{-xr_v}\prod_{u\to v}(u/v)^{-x},
\]
which is the product in~\eqref{eq:mixed}. Vertex $1$ has $\psi_1=1$ and
label $1$, so it introduces no beta factor or label power.
Fatou's lemma passes to the limiting weights.

For the conditional assertion take $n\ge K$ as well. Only the $K$th
factor changes, to $z^{\ell_K}(1-z)^{h_K}$.
The beta integral~\eqref{eq:beta-factor} has the uniform lower bound
$C^{-1}K^{-\ell_K}$ at that site, for fixed exponent bounds.
Its replacement therefore costs at most
$C K^{\ell_K}(1-z)^{h_K}$. Fatou's lemma applies with $z$ fixed,
proving the stated conditional estimate.
\end{proof}

\subsection{Summing the labels}
The following elementary observation is what makes early bases dominate.
It will also control unions of overlapping bases.

\begin{lemma}\label{lem:sums}
Let $F$ be an ordered simple graph with minimum degree at least $d\ge2$,
and let $r_v\ge1$ be integers. Put
\[
 \gamma_v=-(1-x)(r_v+d_-(v))-x d_+(v).
\]
If $0<x<1/2$, the sum of $\prod_vv^{\gamma_v}$ over all increasing
positive vertex labels is finite. For any fixed ordered exponent sequence
$\gamma_1,\ldots,\gamma_h$ obtained this way,
\begin{align}
 \sum_{K<v_1<\cdots<v_h}\prod_i v_i^{\gamma_i}
 &\le C K^{h+\sum_i\gamma_i},\label{eq:anchored-all}\\
 \sum_{K=v_1<v_2<\cdots<v_h}\prod_i v_i^{\gamma_i}
 &\le C K^{h-1+\sum_i\gamma_i}.\label{eq:anchored-one}
\end{align}
\end{lemma}
\begin{proof}
For a nonempty terminal segment $U$ of the ordered vertices, all edges
crossing its boundary point toward earlier vertices outside $U$. Hence
\begin{align}
 \sum_{v\in U}(\gamma_v+1)
 &=|U|-(1-x)\sum_{v\in U}r_v-|E(U)|-x|E(U,U^c)|\notag\\
 &\le x|U|-|E(U)|-x|E(U,U^c)|\notag\\
 &\le-(d-1)x|U|-(1-2x)|E(U)|<0.\label{eq:suffix}
\end{align}
The last step uses $2|E(U)|+|E(U,U^c)|\ge d|U|$.
Place $v_i$ in the dyadic interval $[2^{n_i},2^{n_i+1})$.
The sum in those intervals is at most a constant times
$2^{\sum_i(\gamma_i+1)n_i}$, with $n_1\le\cdots\le n_h$.
Writing the $n_i$ as successive nonnegative differences turns the full sum
into a product of geometric sums, whose exponents are precisely the
negative suffix sums~\eqref{eq:suffix}.
For~\eqref{eq:anchored-all}, start the dyadic intervals at $K$; their
initial factor is $K^{h+\sum_i\gamma_i}$.
For~\eqref{eq:anchored-one}, fix the first label and sum the remaining
$h-1$ labels in the same way. All their suffix sums are still negative,
and the initial factor is $K^{h-1+\sum_i\gamma_i}$.
\end{proof}

For an octahedral graph $F$ on ordered labels $v_1<\cdots<v_k$, write
\begin{equation}\label{eq:f-base}
 f_F(v_1,\ldots,v_k)
 =\prod_{i=1}^k v_i^{-(1-x)(1+d_-(v_i))-x d_+(v_i)}.
\end{equation}
There are finitely many perfect matchings on $k$ ordered positions.
Lemma~\ref{lem:mixed} and Lemma~\ref{lem:sums}, with $d=k-2\ge2$,
give a summable bound on $\E[I_S\prod_{v\in S}D_v]$.
In particular $\E A<\infty$.

\begin{proposition}\label{prop:moments}
The series $A$ has every positive moment.
\end{proposition}
\begin{proof}
For an integer $p\ge1$, expand the $p$th power of finite partial sums
of~\eqref{eq:coefficient}. Classify the $p$ bases by their overlap pattern
and by their missing matchings. There are finitely many such patterns at
fixed $p$ and $k$. In each pattern, the union of the required edges forms
a graph $F$ of minimum degree at least $k-2$, because every vertex belongs
to an octahedral base. Its weight multiplicity $r_v$ is the number of bases
containing it, so $1\le r_v\le p$.
Ignoring extra nonedge requirements only increases the probability; an
inconsistent pattern or a required outdegree larger than $m$ contributes
zero. Otherwise~\eqref{eq:mixed} bounds the weighted expectation, and
Lemma~\ref{lem:sums} sums its labels. This bound is finite for every
pattern. Monotone convergence proves $\E A^p<\infty$.
For an arbitrary real $p>0$, use $A^p\le1+A^{\lceil p\rceil}$.
\end{proof}

\subsection{Common neighbors and the complete count}
For a fixed $k$-set $S$ and $t>\max S$, let $Y_{S,t}$ indicate that $t$
is joined to every vertex of $S$. Put
$N_S(T)=\sum_{\max S<t\le T}Y_{S,t}$.

\begin{lemma}\label{lem:neighbors}
For every fixed $S$,
\begin{equation}\label{eq:fixed-count}
 \frac{N_S(T)}{T^\rho}
 \longrightarrow \frac{\fall{m}{k}}\rho\prod_{v\in S}D_v
 \quad\text{almost surely and in }L^1.
\end{equation}
Moreover $\sup_T\E[(N_S(T)/T^\rho)^2]<\infty$.
\end{lemma}
\begin{proof}
Condition on the whole environment. Distinct apex labels use independent
edge uniforms. There are $\fall{m}{k}$ injective assignments of edge slots
to the vertices of $S$. A union bound and the two-term Bonferroni inequality
give
\begin{equation}\label{eq:occupancy}
 \Prob(Y_{S,t}=1\mid(\psi_j))
 =\fall{m}{k}\prod_{v\in S}p(v,t)
 +O_m\left(\max_{v\in S}p(v,t)\prod_{v\in S}p(v,t)\right).
\end{equation}
Indeed, two distinct compatible assignments specify at least one extra
edge; incompatible assignments have zero intersection probability.
Since $t^xp(v,t)\to D_v$, the conditional mean of $N_S(T)$ is
asymptotic to $\fall{m}{k}\prod D_v\,T^\rho/\rho$.
This follows by bounding the probabilities above and below by constant
multiples of $t^{-kx}$ after any fixed accuracy threshold and summing
$\sum_{t\le T}t^{-kx}\sim T^\rho/\rho$.

The centered variables
$(Y_{S,t}-\E[Y_{S,t}\mid(\psi_j)])/t^\rho$ are independent for this fixed
environment, and their variance sum is bounded, apart from a finite prefix,
by $C(\psi)\sum_t t^{kx-2}<\infty$.
The independent series criterion and Kronecker's lemma
\cite[Theorems~2.5.6 and~2.5.9]{durrett} show that the centered partial
sum is $o(T^\rho)$ almost surely. Fubini transfers this conclusion to the
joint probability space.

For the second moment, two distinct apices $t,s>\max S$ joined to the
same fixed set have probability at most $C_S t^{-kx}s^{-kx}$ by
Lemma~\ref{lem:mixed} with zero weight multiplicities. The diagonal
probabilities are at most $C_S t^{-kx}$. Summing proves the uniform
second-moment assertion. An almost-surely convergent family $Z_T$ with
bounded second moments converges in $L^1$: Fatou bounds the second moment
of its limit $Z$, and Cauchy--Schwarz gives
\[
 \E|Z_T-Z|\le\varepsilon+
 (\E|Z_T-Z|^2)^{1/2}\Prob(|Z_T-Z|>\varepsilon)^{1/2}.
\]
The probability tends to zero and then $\varepsilon$ can decrease to zero.
This proves~\eqref{eq:fixed-count}.
\end{proof}

Let $C_t$ count the induced octahedral $k$-sets in the graph on the
distinct older neighbors of $t$, and let $C(T)=\sum_{3\le t\le T}C_t$.
These counts use only finite graphs, and
$C(T)=\sum_S I_S N_S(T)$.
For a cutoff $K$, let $C^K(T)$ retain bases $S\subset[K]$ and apices $t>K$,
and let $A_K=\sum_{S\subset[K],\,|S|=k}I_S\prod_{v\in S}D_v$.

\begin{proposition}\label{prop:count}
The complete count satisfies
\begin{equation}\label{eq:count-limit}
 \frac{C(T)}{T^\rho}\xrightarrow{L^1}\frac{\fall{m}{k}}\rho A.
\end{equation}
\end{proposition}
\begin{proof}
At fixed $K$, Lemma~\ref{lem:neighbors}, multiplied by each bounded
$I_S$, proves $C^K(T)/T^\rho\to\fall{m}{k}A_K/\rho$ in $L^1$.
No independence between $I_S$ and the weights is required.
The required base edges and the $k$ apex edges give
\[
 \Prob(I_S=1,Y_{S,t}=1)\le C t^{-kx} f_F(S),
\]
summed over the finitely many matchings $F$ on $S$. Therefore
\begin{equation}\label{eq:cutoff-count}
 \limsup_{T\to\infty}\E\frac{C(T)-C^K(T)}{T^\rho}
 \le C\sum_F\sum_{\max S>K} f_F(S).
\end{equation}
The omitted apices $t\le K$ have count at most
$K\binom{m}{k}$, which disappears before the displayed limsup.
The right side tends to zero by Lemma~\ref{lem:sums}, and
$A_K\uparrow A$ in $L^1$. Taking $T\to\infty$ first and then
$K\to\infty$ proves~\eqref{eq:count-limit}.
\end{proof}

\section{Finite homology and link cores}\label{sec:homology}
For vertex $t$, let $L_t$ be its link in $X_t$, the clique complex of the
simple graph on its older neighbors. It has at most $m$ vertices.
The inclusions $L_t\subset X_{t-1}$ induce maps $f_{j,t}$ on $H_j$.
The decomposition $X_t=X_{t-1}\cup\operatorname{St}_{X_t}(t)$ gives
\begin{equation}\label{eq:rank-increment}
 \beta_q(t)-\beta_q(t-1)
 =\dim\ker f_{q-1,t}-\rank f_{q,t}.
\end{equation}
This is the link-rank identity used in~\cite[Proposition~14 of the
author preprint]{ssyh}.
For completeness, the rational simplicial-chain sequence for a union of
subcomplexes $\mathcal K=\mathcal U\cup\mathcal V$, $\mathcal L=\mathcal U\cap\mathcal V$, is
\[
 0\longrightarrow C_*(\mathcal L)\xrightarrow{\zeta\mapsto(\zeta,-\zeta)}
 C_*(\mathcal U)\oplus C_*(\mathcal V)\xrightarrow{(\xi,\eta)\mapsto\xi+\eta}
 C_*(\mathcal K)\longrightarrow0.
\]
It is exact degree by degree because the simplex bases shared by $\mathcal U$ and
$\mathcal V$ are exactly those of $\mathcal L$. If a cycle in $\mathcal K$ is lifted to $(\xi,\eta)$,
its boundary has the form $(\zeta,-\zeta)$, with $\zeta$ a $\mathcal L$-cycle. The class of
$\zeta$ is unchanged by a different lift or a homologous cycle. This gives the
connecting map; a zero connecting class is exactly a lift that can be
modified into two cycles. The other two exact positions follow by lifting
a bounding chain and subtracting it. Thus the sequence yields the homology
exact sequence. A cone has zero reduced homology: adjoining its apex to
an augmented simplex gives a chain homotopy $h$ with $dh+hd=\mathrm{id}$.
Applying the sequence with $\mathcal V$ the star of $t$ gives the exact segment
\[
 H_q(L_t)\xrightarrow{f_{q,t}}H_q(X_{t-1})\longrightarrow H_q(X_t)
 \longrightarrow H_{q-1}(L_t)\xrightarrow{f_{q-1,t}}H_{q-1}(X_{t-1}).
\]
Exactness at the middle terms gives
\[
 \dim H_q(X_t)=\beta_q(t-1)-\rank f_{q,t}+\dim\ker f_{q-1,t},
\]
which is~\eqref{eq:rank-increment}. The same calculation will keep track of the inclusion
maps when vertices are deleted from a link.

\subsection{The typical link}
Put $d=k-2=2q-2$. Delete, repeatedly, vertices of a link graph whose
current degree is less than $d$. The remaining graph is its $d$-core.
Call $L_t$ exceptional if this core has at least $k+1$ vertices.
Let $B(T)$ be the number of exceptional links up to time $T$.
In homology expressions, an octahedral base $S$ denotes its induced clique
complex, rather than only its vertex set.

\begin{lemma}\label{lem:core}
If $L_t$ is not exceptional, then
\[
 \dim H_{q-1}(L_t)=C_t\in\{0,1\},\qquad H_q(L_t)=0.
\]
When $C_t=1$, its unique octahedral base $S$ induces an isomorphism
$H_{q-1}(S)\to H_{q-1}(L_t)$.
\end{lemma}
\begin{proof}
A flag complex with nonzero reduced $H_j$ has at least $2j+2$ vertices.
For $j=0$ this is disconnectedness. For $j\ge1$, take an induced subcomplex
minimal by vertices with nonzero $H_j$. Deleting any vertex makes $H_j$
zero, so the homology exact sequence injects $H_j$ into reduced
$H_{j-1}$ of its link. Induction gives at least $2j$ neighbors for each
vertex. A universal vertex would make the complex a cone; hence there is
also a nonneighbor, proving the lower bound $2j+2$.

A vertex of degree below $2q-2$ has a link on at most $2q-3$ vertices.
Its reduced $H_{q-2}$ and $H_{q-1}$ vanish by the support bound. The exact
sequence therefore shows that its deletion induces an isomorphism on
$H_{q-1}$. It also induces an isomorphism on $H_q$, by the same bound
in the next degree. Every induced octahedral base survives the deletions,
since all its internal degrees are $k-2$.

If the remaining core has fewer than $k$ vertices, both the relevant
homology and its octahedral count vanish. With exactly $k$ vertices,
either there is a universal vertex, and the complex is a cone, or all
degrees are $k-2$. In the latter case its complement is a perfect matching,
so it is the unique octahedral base. Its clique complex is a join of $q$
two-point spaces: splitting at one missing pair expresses it as two cones
over the $(q-1)$-fold join. The reduced exact sequence shifts homology by
one. Starting with two points shows that $H_{q-1}$ is one-dimensional
and $H_q=0$. The isomorphisms under deletion give the asserted inclusion map.
\end{proof}

\begin{lemma}\label{lem:exceptional}
The exceptional links have $\E B(T)=o(T^\rho)$.
\end{lemma}
\begin{proof}
An exceptional link contains a graph $F$ on $h$ vertices of minimum degree
at least $d$, where $k+1\le h\le m$, and $t$ is joined to all these
vertices. There are finitely many ordered graph patterns for fixed $m$.
The $h$ apex edges add one incoming edge to each base vertex.
Lemma~\ref{lem:mixed}, with zero weight multiplicities, therefore bounds
the probability of a specified occurrence by
\[
 C t^{-hx}\prod_{v\in F}v^{-(1-x)(1+d_-(v))-x d_+(v)}.
\]
Lemma~\ref{lem:sums} sums the base labels. Summing the patterns gives
$\Prob(L_t\text{ exceptional})\le C t^{-(k+1)x}$.
Thus $\E B(T)$ is $O(T^{1-(k+1)x})$, $O(\log T)$ or $O(1)$ according
as $1-(k+1)x$ is positive, zero or negative. Each is $o(T^{1-kx})$.
For $m=k$, no exceptional link is possible.
\end{proof}

\subsection{Earlier cones and the Betti limit}
\begin{lemma}\label{lem:cones}
Almost surely $A>0$. For every $K$, each existing octahedral base
$S\subset[K]$ has a common neighbor $s_S>K$ at a finite random time.
The maximum of their finitely many first such times is finite almost surely.
\end{lemma}
\begin{proof}
The occupancy calculation~\eqref{eq:occupancy} applies to any fixed
$h$-set with $h\le k$, and its conditional common-neighbor probability
is asymptotic to a positive multiple of $t^{-hx}$.
It has divergent sum because $hx<1$. Conditional independence and the
second Borel--Cantelli lemma~\cite[Theorem~2.3.7]{durrett} give infinitely
many common neighbors. If one fixed extra vertex is forbidden, the
probability of also joining it is $O(t^{-(h+1)x})$, whenever $h+1\le m$,
so there are still infinitely many common neighbors avoiding it.
These statements hold simultaneously for all fixed finite sets and extra
vertices, by countability and Fubini.

A nonadjacent pair exists once a new vertex has more than $m$ old vertices
available. Starting from this two-point sphere, choose a common neighbor
of the current octahedral sphere, then a second common neighbor avoiding
the first. The two new vertices give its suspension as an induced
subgraph. Repeat until there are $k$ vertices. At each step the current
sphere has at most $k-2$ vertices, so the required edges and the avoidance
condition are within $m$. Edges between old vertices are never added later.
This constructs an induced octahedral base. Its weights are positive,
hence $A>0$.

For a fixed $K$ there are finitely many bases in $[K]$. Each has a common
neighbor $s>K$ almost surely. Let $\tau_K$ be the largest of their first
such labels, or $K$ if no base exists. Then $\tau_K<\infty$ almost surely.
The star of each such $s$ contains the cone on its base, so the map from
that base's homology into $H_{q-1}(X_{t-1})$ is zero for $t>\tau_K$.
\end{proof}

\begin{proposition}\label{prop:betti}
The limits in~\eqref{eq:betti-limit} hold.
\end{proposition}
\begin{proof}
On a nonexceptional link after $\tau_K$, if its unique base is in $[K]$,
Lemma~\ref{lem:core} identifies its entire $H_{q-1}$ with that of the
base. The latter maps to zero in the old complex by Lemma~\ref{lem:cones}.
Equation~\eqref{eq:rank-increment} therefore gives a birth of dimension one,
and there is no death because $H_q(L_t)=0$. All other nonexceptional links
give nonnegative increments. On exceptional links both homology dimensions
and base counts have a bound depending only on $m,q$, since the link has
at most $m$ vertices. Summing~\eqref{eq:rank-increment} gives, for a fixed
constant $M$,
\begin{align}
 \beta_q(T)&\le C(T)+M B(T),\label{eq:hom-upper}\\
 \beta_q(T)&\ge C^K(T)-C^K(\min(T,\tau_K))-M B(T).
 \label{eq:hom-lower}
\end{align}
The prefix term divided by $T^\rho$ tends to zero almost surely, because
$\tau_K$ is finite. It tends to zero in $L^1$ as well: it is bounded by
$C^K(T)/T^\rho$, and the latter has uniformly bounded second moments by
Lemma~\ref{lem:neighbors} and the finite number of bases in $[K]$.
The Cauchy--Schwarz estimate in that lemma applies. No expectation bound
on $\tau_K$ is needed.

It follows from~\eqref{eq:hom-upper}--\eqref{eq:hom-lower} that
\[
 |\beta_q(T)-C(T)|
 \le C(T)-C^K(T)+C^K(\min(T,\tau_K))+M B(T).
\]
Use~\eqref{eq:cutoff-count} and Lemma~\ref{lem:exceptional}, taking
$T\to\infty$ at fixed $K$ and then $K\to\infty$. This difference is
$o(T^\rho)$ in $L^1$. Proposition~\ref{prop:count} proves the first
limit in~\eqref{eq:betti-limit}; taking expectations proves the second.
Its limiting constant is positive by Lemma~\ref{lem:cones}.
Dividing by that deterministic mean, which is positive for sufficiently
large $T$, yields $L^1$ convergence of the
mean-normalized variables to $A/\E A$ in this coupling, and hence the
asserted convergence in distribution.
\end{proof}

\section{Nondegeneracy and the upper tail}\label{sec:tail}
The coefficient retains random early structure. Its nonconstancy can be
seen by making the weights of every vertex before a large fixed label
small at the same time.

\begin{proposition}\label{prop:support}
The coefficient $A$ has essential infimum zero.
\end{proposition}
\begin{proof}
Let $d=k-2$ and $e=k(q-1)$, the degree and edge count of an octahedral
base. Condition on $\psi_K=z$, for $K\ge2$, and split the series into
bases with a label below $K$, bases with minimum $K$, and bases entirely
after $K$.

Recall that $\ell_K=r_K+d_-(K)$, while $h_K$ counts weights below $K$
and required edges crossing $K$ in Lemma~\ref{lem:mixed}. For a single base,
$r_K$ is one if $K$ belongs to it and zero otherwise.
In the first group, at least one weight contains $1-\psi_K$, so $h_K\ge1$.
Also $\ell_K\le d+1$. Summing its conditional bound using the same
summable base weights~\eqref{eq:f-base} gives
\begin{equation}\label{eq:old-support}
 \E[A_{\mathrm{old}}\mid\psi_K=z]
 \le C(1+K^{d+1})(1-z).
\end{equation}
For a base with minimum exactly $K$, the vertex $K$ has weight
multiplicity one, incoming degree $d$ and outgoing degree zero. The
conditional cost is at most $C K^{d+1}$.
The sum of all base exponents in~\eqref{eq:f-base} is $-(1-x)k-e$.
Using~\eqref{eq:anchored-one},
\begin{align}
 \E[A_{\min=K}\mid\psi_K=z]
 &\le C K^{d+1}K^{k-1-(1-x)k-e}\notag\\
 &=C K^{kx-e+d}=C K^{kx-2(q-1)^2}.\label{eq:min-support}
\end{align}
Bases entirely after $K$, including their weights, are independent of
$\psi_K$. Equation~\eqref{eq:anchored-all} gives
\begin{equation}\label{eq:after-support}
 \E[A_{\min>K}\mid\psi_K=z]\le C K^{kx-e}.
\end{equation}
Both exponents in~\eqref{eq:min-support}--\eqref{eq:after-support} are
negative, since $kx<1$, $e-d=2(q-1)^2\ge2$ and $e\ge4$.
Given $\varepsilon>0$, first choose $K$ so that their sum is below
$\varepsilon/2$, and then $\eta\in(0,1)$ so that the right side
of~\eqref{eq:old-support} is below $\varepsilon/2$ when
$1-\eta<z<1$. This Beta interval has positive probability.
The conditional mean of $A$ on the interval is below $\varepsilon$;
Markov's inequality therefore gives
$\Prob(A\le2\varepsilon)>0$. This proves $\essinf A=0$.
\end{proof}

Together with $A>0$ almost surely and $0<\E A<\infty$, this proves that
$W=A/\E A$ is nondegenerate.
Proposition~\ref{prop:moments} gives every positive moment of $W$.
For every $p>0$ its survival function is $O_p(z^{-p})$, by Markov's
inequality. To finish Theorem~\ref{thm:main}, suppose~\eqref{eq:rv} held
with a finite $\alpha>0$. Its survival function must be positive
eventually. For any $\epsilon>0$, the ratio at $u=2$ is at least
$2^{-(\alpha+\epsilon)}$ for large $z$. Iteration gives a dyadic lower
bound, and monotonicity extends it to
$\Prob(W>y)\ge C y^{-(\alpha+\epsilon)}$ for all sufficiently large $y$.
An integer $p>\alpha+\epsilon$ then satisfies
\[
 \E W^p=p\int_0^\infty y^{p-1}\Prob(W>y)\,dy=\infty,
\]
by nonnegative Fubini, contradicting Proposition~\ref{prop:moments}.
A bounded survival support also cannot satisfy~\eqref{eq:rv}.
Thus no finite positive regular-variation exponent occurs in the strict
regime of the theorem.

\section*{AI-use statement}
OpenAI Codex was used extensively for literature review, development and
checking of the arguments, and preparation of this manuscript. The
mathematical claims are offered for independent review.

\begin{thebibliography}{9}
\bibitem{dey-terlov}
P.~S.~Dey and G.~Terlov,
\emph{Limiting distributions of triangle counts in linear preferential attachment models},
preprint, arXiv:2605.28776v1, 2026.
\url{https://arxiv.org/abs/2605.28776v1}.

\bibitem{durrett}
R.~Durrett, \emph{Probability: Theory and Examples}, fifth edition,
Cambridge University Press, 2019.
Author draft, Version January 11, 2019:
\url{https://services.math.duke.edu/~rtd/PTE/PTE5_011119.pdf}.

\bibitem{siu}
C.~Siu, \emph{The Topological Behavior of Preferential Attachment Graphs},
SIAM Journal on Applied Algebra and Geometry \textbf{9}(2) (2025), 432--455.
\href{https://doi.org/10.1137/24M1693179}{doi:10.1137/24M1693179}.
Author preprint: \url{https://arxiv.org/abs/2406.17619v2}.

\bibitem{ssyh}
C.~Siu, G.~Samorodnitsky, C.~L.~Yu and R.~He,
\emph{The asymptotics of the expected Betti numbers of preferential attachment clique complexes},
Advances in Applied Probability \textbf{57}(3) (2025), 940--968.
\href{https://doi.org/10.1017/apr.2024.66}{doi:10.1017/apr.2024.66}.
Author preprint: \url{https://arxiv.org/abs/2305.11259v3}.
\end{thebibliography}
\end{document}
